\documentclass[12pt]{article}
\usepackage[utf8]{inputenc}
\usepackage[T1]{fontenc}
\usepackage{hyperref}
\title{Pam Language Resources for K-12 Curriculum Development}
\author{Compiled by Claude Sonnet 5.5 (claude-sonnet-5-5) for Andr\'as Kornai}
\date{}
\begin{document}
\maketitle

\section{Introduction}
Pam (ISO 639-3 pmn; Glottocode pamm1242) is a nearly extinct language of northern Cameroon, spoken around Tcholliré, with only a few dozen speakers reported and a shift to Adamawa Fulfulde. Search found no teaching materials; only reference pages.

\section{Primary Educational Resources}
None found.

\section{Linguistic and Reference Materials}
\begin{itemize}
\item Wikipedia, Pam language: \url{https://en.wikipedia.org/wiki/Pam_language}
\item Glottolog, Pam (pamm1242): \url{https://glottolog.org/resource/languoid/id/pamm1242.bigmap.html}
\item Numeral data page for Pam, Max Planck EVA: \url{https://lingweb.eva.mpg.de/channumerals/Pam.htm}
\end{itemize}

\section{Religious and Community Resources}
None found.

\section{Digital and Audio Resources}
None found.

\section{Curriculum Development Context}
With very few remaining speakers and no known documentation beyond word lists such as numerals, any heritage programme would depend on fresh documentation work with the remaining speakers and community.
\end{document}